\chapter{Giải pháp và công nghệ sử dụng}
\label{chap:solution-tech}

Để xây dựng một nền tảng IoT hoàn chỉnh vừa đảm bảo khả năng phản ứng an toàn tại phần cứng, vừa đủ gọn nhẹ để vận hành ổn định trên một máy chủ ảo chi phí thấp, nhóm chúng em cần chọn lọc kỹ từng công nghệ ở mỗi tầng. Thay vì ghép nối nhiều dịch vụ phức tạp, tiêu chí xuyên suốt của nhóm là ưu tiên các công cụ mã nguồn mở ổn định, có tài liệu rõ ràng, dễ kiểm thử tự động và tận dụng tối đa tài nguyên sẵn có. Chương này trình bày tổng quan giải pháp 4 tầng, vai trò cụ thể của từng công nghệ, lý do cho các quyết định kiến trúc ngay trong bối cảnh từng tầng kỹ thuật, cũng như cơ sở chuẩn hóa cho dữ liệu học máy.

\section{Tổng quan giải pháp kỹ thuật 4 tầng End-to-End}
\label{sec:c2-overview}

Hệ thống GreenHub được chia thành bốn tầng chức năng tách bạch rõ trách nhiệm, giúp luồng dữ liệu từ cảm biến đến người dùng cuối đi theo một trục thống nhất và dễ cô lập sự cố khi vận hành:

\begin{enumerate}[leftmargin=*]
    \item \textbf{Tầng thiết bị biên (Edge / Device Layer):} Gồm các nút cảm biến (Sensor Node) đọc thông số môi trường theo chu kỳ cố định (\(5\;\text{s}\) ở chế độ \texttt{DEMO\_MODE} hoặc cấu hình từ xa lên tới \(60\;\text{s}\) hỗ trợ ngủ sâu \texttt{deep sleep}), đóng gói thành khung dữ liệu nhị phân chỉ \(20\;\text{byte}\) và truyền không dây qua ESP-NOW tới Gateway. Gateway ESP32 đóng vai trò cầu nối ESP-NOW \(\leftrightarrow\) MQTT over TLS cổng \texttt{8883}, đồng thời gánh vác khâu tiền xử lý nhiễu (loại gai đơn lẻ, lọc trung vị 3 mẫu kết hợp EMA), lưu đệm 150 bản tin trên RAM chống mất dữ liệu khi đứt mạng và thực thi lớp bảo vệ phần cứng cùng failsafe cục bộ cho 3 kênh relay (\texttt{pump\_1}, \texttt{pump\_2}, \texttt{fan}).
    \item \textbf{Tầng kết nối và bảo mật truyền thông (Connectivity Layer):} Đặt trọng tâm vào broker Eclipse Mosquitto 2.x cấu hình bắt buộc mã hóa TLS 1.2+ trên cổng \texttt{8883}, vô hiệu hóa truy cập ẩn danh và quản lý tài khoản cùng quyền truy cập động qua Dynamic Security Plugin. Mỗi gateway chỉ được cấp quyền đọc và ghi trong đúng nhánh topic của nhà kính mình (\texttt{gh/\{farm\_id\}/\{greenhouse\_id\}/}).
    \item \textbf{Tầng nền tảng lõi (Core Platform Layer):} Thiết kế theo dạng monolith hai tiến trình chạy độc lập từ chung một mã nguồn Python 3.12. Tiến trình \texttt{api} (FastAPI + Uvicorn) phụ trách REST API, xác thực JWT kết hợp băm mật khẩu Argon2id (một quyền người dùng duy nhất \texttt{owner}), cô lập đa tenant theo \texttt{farm\_id}, cấp phát gateway và đẩy sự kiện thời gian thực qua WebSocket. Tiến trình \texttt{worker} (\texttt{aiomqtt} + APScheduler) lắng nghe toàn bộ nhánh \texttt{gh/\#} để thực hiện chuỗi xử lý: thu nhận dữ liệu, kiểm định chất lượng, tính VPD/DLI, chạy bộ luật điều khiển, quản lý vòng đời lệnh -- cảnh báo -- thông báo và chạy suy luận học máy định kỳ. Cơ sở dữ liệu PostgreSQL 16 tích hợp TimescaleDB 2.17 vừa lưu trữ dữ liệu quan hệ lẫn chuỗi thời gian, vừa đảm nhiệm vai trò kênh truyền tin nội bộ tốc độ cao giữa \texttt{api} và \texttt{worker} thông qua \texttt{LISTEN/NOTIFY}.
    \item \textbf{Tầng ứng dụng, học máy và hạ tầng triển khai (Application, AI \& DevOps Layer):} Cung cấp giao diện Web Console PWA viết bằng React 18, TypeScript và Vite gồm 9 trang chức năng tiếng Việt, hỗ trợ cập nhật tức thời và hiển thị cảnh báo thời gian thực qua WebSocket. Quy trình huấn luyện LightGBM chạy ngoại tuyến qua các kịch bản CLI trong thư mục \texttt{ml/}, đi qua cổng kiểm định chất lượng rồi xuất tệp mô hình \texttt{.joblib} cho \texttt{worker} nạp vào suy luận. Toàn bộ hệ thống được đóng gói thành 5 container Docker Compose đứng sau Caddy 2.8 tự động cấp phát HTTPS Let's Encrypt và chia sẻ chứng chỉ TLS cho Mosquitto.
\end{enumerate}

\section{Tầng thiết bị biên và lập trình nhúng}
\label{sec:c2-edge}

Ở tầng thiết bị biên, nhóm lựa chọn bộ công cụ \textbf{PlatformIO} kết hợp khung lập trình Arduino cho hai dòng vi điều khiển phổ biến là \textbf{ESP32} (\texttt{espressif32@\textasciicircum 6.9.0}) và \textbf{ESP8266} (\texttt{espressif8266@\textasciicircum 4.2.1})~\cite{platformioDocs}. PlatformIO giúp quản lý thư viện phụ thuộc theo phiên bản ghim cố định và tổ chức nhiều cấu hình biên dịch (\texttt{env}) trong cùng một mã nguồn, nhờ đó một bản firmware duy nhất cho Sensor Node có thể biên dịch linh hoạt cho nhiều biến thể phần cứng thực tế như AHT20, SHT30, DHT11 hay bo mạch ESP8266 chỉ lắp cảm biến ẩm đất.

Đối với liên kết không dây từ các nút cảm biến về Gateway, nhóm quyết định sử dụng giao thức \textbf{ESP-NOW} của Espressif thay vì cho từng Sensor Node kết nối Wi-Fi trực tiếp tới broker đám mây~\cite{espressifEspNowDocs}. Do ESP-NOW truyền dữ liệu trực tiếp giữa các địa chỉ MAC ở tầng liên kết dữ liệu trên băng tần \(2{,}4\;\text{GHz}\) (tải trọng tối đa \(250\;\text{byte}\)), mỗi nút cảm biến chỉ mất vài mili-giây để phát xong gói tin nhị phân \(20\;\text{byte}\), loại bỏ hoàn toàn chi phí bắt tay DHCP, TCP và TLS lặp lại sau mỗi chu kỳ ngủ sâu (\texttt{deep sleep}). Bên cạnh lợi thế tiết kiệm năng lượng, kiến trúc qua trung gian này còn gom toàn bộ chứng chỉ TLS, hàng đợi RAM 150 bản tin chống mất mạng và cơ chế an toàn/failsafe phần cứng về duy nhất một điểm tại Gateway, giúp việc quản trị, đồng bộ cấu hình (\texttt{cfg\_version}) và bảo mật thiết bị biên trở nên tập trung và gọn nhẹ hơn nhiều. Các cảm biến môi trường trên nút được giao tiếp qua những thư viện chuẩn ghim phiên bản cố định gồm \texttt{Adafruit AHTX0} (\texttt{@\textasciicircum 2.0.5}), \texttt{Adafruit SHT31} (\texttt{@\textasciicircum 2.2.2}), \texttt{DHT sensor} (\texttt{@\textasciicircum 1.4.6}) và \texttt{BH1750} (\texttt{@\textasciicircum 1.3.0}), trong khi độ ẩm đất điện dung v1.2 được đọc qua bộ chuyển đổi ADC phần cứng kèm chân cấp nguồn điều khiển đóng ngắt trong \(50\;\text{ms}\) nhằm hạn chế ăn mòn điện phân đầu dò.

Tại Gateway ESP32, nhóm tích hợp bộ ba thư viện gồm \texttt{WiFiManager} (\texttt{@\textasciicircum 2.0.17}) để tạo cổng cấu hình \textit{captive portal} khi khởi chạy lần đầu hoặc khi nhấn giữ nút \texttt{BOOT} nhằm nạp thông số Wi-Fi và MQTT vào vùng nhớ NVS, \texttt{WiFiClientSecure} kết hợp \texttt{PubSubClient} (\texttt{@\textasciicircum 2.8}) để duy trì phiên MQTT over TLS có xác thực chứng chỉ gốc CA, cùng \texttt{ArduinoJson} (\texttt{@\textasciicircum 7.2.1}) để đóng gói và giải mã bản tin JSON gọn nhẹ, hạn chế phân mảnh bộ nhớ heap trên vi điều khiển.

\section{Tầng kết nối và bảo mật truyền thông}
\label{sec:c2-connectivity}

Cầu nối trung tâm giữa thiết bị biên và nền tảng máy chủ là \textbf{Eclipse Mosquitto} phiên bản \texttt{2.0.20}~\cite{mosquittoDocs}. Mosquitto nổi bật ở mức tiêu thụ bộ nhớ rất thấp (giới hạn \(128\;\text{MB}\) RAM vẫn vận hành dư dả hàng trăm kết nối đồng thời), hỗ trợ đầy đủ các cơ chế thiết yếu cho IoT như đảm bảo chuyển phát \texttt{QoS 1}, thông điệp lưu giữ (\texttt{retained}) cho trạng thái và cấu hình, cùng cơ chế di chúc \texttt{LWT} giúp máy chủ nhận diện tức thì khi một gateway bị mất nguồn hoặc rớt mạng đột ngột.

Điểm đáng chú ý trong cách triển khai của nhóm là việc bật \textbf{Dynamic Security Plugin} của Mosquitto. Thay vì phải chỉnh tay tệp mật khẩu và nạp lại tiến trình broker mỗi lần thêm nhà kính mới, dịch vụ \texttt{api} có thể gửi lệnh quản trị qua topic hệ thống \texttt{\$CONTROL/dynamic-security/v1} để tạo ngay tài khoản MQTT và gắn bộ luật ACL cô lập hoàn toàn từng gateway ngay tại thời điểm cấp phát thiết bị. Đồng thời, cổng công khai \texttt{8883} bắt buộc mã hóa TLS 1.2+ sử dụng trực tiếp cặp chứng chỉ Let's Encrypt chia sẻ từ Caddy (chi tiết cơ chế đồng bộ tại Mục~\ref{sec:c2-devops}), nhờ đó tránh phải duy trì hai hệ thống cấp phát chứng chỉ song song.

\section{Tầng nền tảng lõi và lưu trữ dữ liệu}
\label{sec:c2-platform}

Toàn bộ tầng nền tảng phía máy chủ được nhóm viết bằng \textbf{Python 3.12} và tổ chức theo kiến trúc \textbf{monolith hai tiến trình (\texttt{api} và \texttt{worker}) dùng chung mã nguồn và Docker image} (\texttt{greenhub-platform:latest}), tận dụng mô hình bất đồng bộ (\texttt{asyncio}) để xử lý đồng thời lượng lớn kết nối mạng trong mức tài nguyên khiêm tốn. Việc tách riêng tiến trình \texttt{worker} khỏi \texttt{api} là quyết định kiến trúc nhằm đảm bảo luồng nhận bản tin MQTT liên tục, bộ luật điều khiển và tác vụ suy luận LightGBM không bao giờ chặn hoặc làm chậm tốc độ phản hồi HTTP và WebSocket của \texttt{api}. Dù dùng chung một image đóng gói, nhóm áp đặt ca kiểm thử tự động \texttt{test\_api\_process\_does\_not\_import\_heavy\_ml\_libraries} để khẳng định tiến trình \texttt{api} không bao giờ nạp các thư viện học máy nặng như \texttt{pandas}, \texttt{scikit-learn} hay \texttt{lightgbm}, giữ mức tiêu thụ RAM của \texttt{api} an toàn dưới ngưỡng \(384\;\text{MB}\).

Cụ thể, tiến trình \texttt{api} sử dụng \textbf{FastAPI} (\texttt{0.162.0}) chạy trên máy chủ ASGI \textbf{Uvicorn} (\texttt{0.54.0}) kết hợp \textbf{Pydantic v2} (\texttt{2.14.2}) để kiểm tra chặt dữ liệu vào/ra và tự sinh tài liệu đặc tả OpenAPI~\cite{fastapiDocs}. Lớp truy cập cơ sở dữ liệu dùng \textbf{SQLAlchemy 2 Async} (\texttt{2.1.8}) đi cùng trình điều khiển hiệu năng cao \texttt{asyncpg} (\texttt{0.32.0}) và công cụ quản lý phiên bản lược đồ \texttt{Alembic} (\texttt{1.20.1})~\cite{sqlalchemyDocs}. Về bảo mật xác thực, mật khẩu người dùng được băm bằng thuật toán kháng tấn công phần cứng \textbf{Argon2id} (\texttt{argon2-cffi==25.1.0})~\cite{argon2Rfc}, còn phiên làm việc trên cả REST API lẫn WebSocket được bảo vệ bằng mã thông báo \textbf{JWT} (\texttt{PyJWT==2.15.1}) gắn chặt với một quyền duy nhất \texttt{owner} và bộ lọc \texttt{farm\_id}. Trong khi đó, tiến trình \texttt{worker} kết hợp thư viện MQTT bất đồng bộ \texttt{aiomqtt} (\texttt{2.5.1}) với bộ lập lịch tác vụ \textbf{APScheduler} (\texttt{3.11.3}) để duy trì song song luồng xử lý sự kiện tức thời từ broker và ba chu kỳ quét định kỳ (\(5\;\text{s}\), \(60\;\text{s}\) và \(10\;\text{phút}\)).

Về lưu trữ và phối hợp liên tiến trình, \textbf{PostgreSQL 16} kết hợp phần mở rộng \textbf{TimescaleDB 2.17.2} giữ vai trò kho dữ liệu duy nhất của toàn hệ thống~\cite{postgresDocs,timescaledbDocs}. PostgreSQL đảm bảo tính toàn vẹn quan hệ cho 16 bảng nghiệp vụ (từ trang trại, nhà kính, thiết bị đến vụ trồng và cảnh báo), trong khi TimescaleDB tự động phân vùng bảng \texttt{readings} thành \textit{Hypertable} theo mốc thời gian và duy trì hai khung nhìn tổng hợp liên tục (\textit{Continuous Aggregates}) là \texttt{readings\_1m} và \texttt{readings\_1h} giúp biểu đồ lịch sử truy vấn nhanh mà không phải quét lại dữ liệu thô. Đặc biệt, thay vì bổ sung thêm một dịch vụ hàng đợi riêng như Redis gây tốn thêm bộ nhớ và tăng thêm điểm lỗi tiềm tàng trên VPS \(2\;\text{GB}\) RAM, nhóm tận dụng trực tiếp cơ chế \textbf{PostgreSQL \texttt{LISTEN/NOTIFY}} qua trình điều khiển duy nhất \texttt{asyncpg} trên hai kênh \texttt{worker\_events} và \texttt{ws\_events}, kết hợp cùng vòng quét bù 5 giây trong \texttt{worker/scheduler.py} để phát lệnh điều khiển và cập nhật WebSocket gần như tức thời.

\section{Tầng ứng dụng Web Console và hệ thống thông báo}
\label{sec:c2-web}

Giao diện quản trị và giám sát của GreenHub được xây dựng theo hướng \textbf{Progressive Web App (PWA)} bằng \textbf{React 18.3.1}, \textbf{TypeScript 5.7} (\texttt{\textasciitilde 5.7.2}) và công cụ đóng gói \textbf{Vite 6.0.5}~\cite{reactDocs,viteDocs}. Lựa chọn kiến trúc PWA cho phép người vận hành sử dụng trực tiếp trên máy tính hoặc cài đặt ra màn hình chính trên điện thoại Android và iOS 16.4+ với trải nghiệm nhất quán từ một mã nguồn web duy nhất, tránh phân mảnh công sức phát triển cho các ứng dụng di động native riêng biệt.

Trong mã nguồn giao diện, nhóm dùng \texttt{react-router-dom} (\texttt{6.28.0}) để tổ chức định tuyến cho 9 trang chức năng tiếng Việt, \texttt{@tanstack/react-query} (\texttt{5.62.0}) để quản lý trạng thái truy vấn và tự động làm mới bộ nhớ đệm khi nhận sự kiện WebSocket, cùng thư viện \texttt{recharts} (\texttt{3.10.1}) để vẽ các biểu đồ vi khí hậu trực quan có tô vùng ngưỡng tối ưu cho VPD, nhiệt độ và ánh sáng.

Đối với cơ chế thông báo sự cố, hệ thống tập trung vào kiến trúc gọn nhẹ gồm hai lớp bổ trợ: thông báo In-app thời gian thực qua WebSocket (\texttt{/ws}) hiển thị Toast và badge chuông trên Web Console khi người dùng trực tuyến, kết hợp với còi Buzzer tại chỗ trên Gateway ESP32 khi thiết bị mất kết nối mạng kéo dài hoặc kích hoạt Failsafe cục bộ, loại bỏ hoàn toàn sự phụ thuộc vào hạ tầng Web Push (VAPID) hay máy chủ gửi thư SMTP bên ngoài.

\section{Phân hệ học máy và xử lý dữ liệu nông học}
\label{sec:c2-ml}

Phân hệ hỗ trợ quyết định canh tác của GreenHub giải quyết duy nhất một bài toán hồi quy toàn cục (\texttt{task = 'yield'}) trên dữ liệu gộp từ nhiều nhà kính: dự đoán năng suất chuẩn hóa (\texttt{relative\_yield}) theo điều kiện khí hậu vụ trồng và dùng chung mô hình đó để xếp hạng gợi ý cây nên trồng. Nhóm lựa chọn \textbf{LightGBM} (\texttt{4.7.0}) làm mô hình hồi quy chính, kết hợp hệ sinh thái \textbf{scikit-learn} (\texttt{1.9.1}), \textbf{pandas} (\texttt{3.1.1}), \textbf{numpy} (\texttt{2.5.4}) và \textbf{joblib} (\texttt{1.6.0}) để tiền xử lý, kiểm định và đóng gói phiên bản mô hình~\cite{lightgbmPaper,scikitLearnDocs}. Quyết định dùng LightGBM xuất phát từ bản chất dữ liệu đầu vào là dữ liệu dạng bảng (16 đặc trưng khí hậu tổng hợp cộng biến phân loại \texttt{crop\_code}, quy mô từ hàng trăm đến vài nghìn mẫu vụ trồng): cơ chế phát triển cây theo lá (\textit{leaf-wise}) của LightGBM nắm bắt tốt các đường đáp ứng sinh học phi tuyến, xử lý tự nhiên biến phân loại \texttt{crop\_code} lẫn giá trị khuyết \texttt{NaN} khi nhà kính không gắn cảm biến sáng BH1750, và suy luận trong vài mili-giây trên CPU mà không cần GPU. Quan trọng hơn cả, phân hệ học máy được tách biệt hoàn toàn khỏi vòng điều khiển đóng ngắt relay, bảo đảm mọi phản ứng phần cứng luôn giữ tính tất định và an toàn tuyệt đối.

Để triệt tiêu hoàn toàn hiện tượng lệch đặc trưng giữa lúc huấn luyện và lúc chạy thật (\textit{training-serving skew}), toàn bộ logic trích xuất 16 đặc trưng khí hậu (\texttt{feature\_version = 1}) được cài đặt duy nhất tại mô-đun \texttt{platform/core/climate\_features.py} bằng \textbf{Python thuần} (không import \texttt{pandas} hay \texttt{numpy}). Nhờ thiết kế này, cả tiến trình \texttt{worker} khi suy luận trực tuyến, pipeline huấn luyện \texttt{ml/} lẫn bộ sinh dữ liệu \texttt{tools/simulator} đều gọi chung một hàm tính toán bất biến, đồng thời bảo vệ gói \texttt{core} và tiến trình \texttt{api} khỏi bị kéo các thư viện ML nặng vào bộ nhớ RAM.

Do số lượng vụ trồng thật trên mô hình vật lý còn hạn chế trong giai đoạn đồ án (một vụ xà lách ngắn ngày cần 3--4 tuần sinh trưởng), bộ dữ liệu mô phỏng ngoại tuyến (\texttt{source = 'synthetic'}) đóng vai trò nguồn huấn luyện chủ lực, được thiết kế chặt chẽ trong \texttt{tools/simulator/growth.py} và \texttt{tools/simulator/crops\_growth.json} theo ba lớp chuẩn mực kỹ thuật để phù hợp thực chất với bài toán nông học -- học máy:

\begin{enumerate}[leftmargin=*]
    \item \textbf{Sinh chuỗi thời gian theo chuẩn vật lý vi khí hậu thay vì lấy mẫu tóm tắt độc lập:} Thay vì bốc ngẫu nhiên trực tiếp 16 con số đặc trưng (dễ tạo ra các tổ hợp phi vật lý giữa nhiệt độ, độ ẩm và VPD), bộ sinh dữ liệu (\texttt{synth-yield}) mô phỏng chuỗi khí hậu theo từng giờ trong cửa sổ chuẩn \(168\;\text{giờ}\) (7 ngày) gắn múi giờ địa phương \texttt{Asia/Ho\_Chi\_Minh}. Trong mỗi ngày, nhiệt độ biến thiên theo sóng hình sin \(24\;\text{giờ}\) kèm dịch pha và nhiễu Gauss; độ ẩm tương đối \(\text{RH}\) dao động ngược pha với nhiệt độ (khô hơn vào ban trưa, ẩm hơn về đêm, chặn trong \([28\%,\; 97\%]\)); bức xạ ánh sáng chỉ xuất hiện trong khung quang kỳ \(06{:}00\text{--}18{:}00\) theo dạng nửa sóng sin nhân hệ số mây ngẫu nhiên (\([0{,}72;\; 1{,}08]\)) và bằng \(0\) vào ban đêm. Tại từng giờ, đại lượng VPD được tính trực tiếp từ \((T, \text{RH})\) theo phương trình áp suất hơi bão hòa chuẩn \textbf{Magnus--Tetens} với độ lệch nhiệt độ lá \(\Delta T_{\text{lá}} = 1{,}0\,^\circ\text{C}\), còn bức xạ quang hợp được quy đổi theo chuẩn ánh sáng tự nhiên \(\text{PPFD} \approx \text{lux} \times 0{,}0185\;\mu\text{mol}\cdot\text{m}^{-2}\cdot\text{s}^{-1}\)~\cite{monteithVpd}. Mô phỏng phân tầng qua \textbf{6 vùng khí hậu đặc trưng của Việt Nam} (\texttt{dalat}, \texttt{hanoi\_winter}, \texttt{hanoi\_summer}, \texttt{hcm}, \texttt{mekong}, \texttt{central\_coast}) kết hợp \textbf{3 mức trang bị nhà kính} (\texttt{basic}, \texttt{standard}, \texttt{full} làm thay đổi độ lệch nhiệt, biên độ ngày -- đêm, độ che nắng và độ ổn định tưới), đồng thời tái tạo hai khuyết tật phổ biến của hệ thống IoT thực tế là tỷ lệ mất gói ngẫu nhiên theo giờ (\([0\%,\; 12\%]\), vẫn đảm bảo \texttt{coverage} \(\ge 0{,}7\)) và khoảng \(5\%\) nhà kính không lắp cảm biến BH1750 (\texttt{dli\_mean = null} trong đặc trưng đầu vào dù sinh trưởng vật lý vẫn phụ thuộc DLI ẩn).

    \item \textbf{Mô hình sinh trưởng cây trồng chuẩn sinh lý thực vật (\textit{Multiplicative Biophysical Response}):} Nhãn năng suất tương đối không phải hàm tuyến tính đơn giản mà tuân theo quy luật tương tác nhân tính giữa các yếu tố giới hạn sinh trưởng trong hàm \texttt{compute\_crop\_relative\_yield}~\cite{monteithVpd,wurChallengeDataset}:
    \begin{equation}
    \text{relative\_yield} = f_T(T_{\text{mean}}) \cdot f_{\text{DLI}}(\text{DLI}) \cdot f_{\text{VPD}}(\text{VPD}_{\text{mean}}) \cdot f_{\text{soil}}(\text{soil}_{\text{mean}}) \cdot f_{\text{heat}} \cdot f_{\text{DIF}} \cdot \text{gh\_effect} \cdot \varepsilon
    \label{eq:synthetic-growth}
    \end{equation}
    Trong đó, các tham số sinh học của từng loại cây (\texttt{tomato\_cherry}, \texttt{strawberry}, \texttt{lettuce}, \texttt{cucumber}, \texttt{bell\_pepper}) được định nghĩa riêng trong \texttt{crops\_growth.json}: các hàm đáp ứng \(f_T, f_{\text{VPD}}, f_{\text{soil}} \in [0, 1]\) mang dạng \textbf{hình thang theo bốn ngưỡng sinh lý} (\(\text{min}, \text{opt}_1, \text{opt}_2, \text{max}\)), đạt giá trị cực đại \(1{,}0\) trong vùng tối ưu và suy giảm tuyến tính ra hai biên chịu đựng; hàm \(f_{\text{DLI}} \in [0, 1]\) mô phỏng đường cong \textbf{bão hòa ánh sáng hàm mũ} chuẩn hóa theo ngưỡng bão hòa từng loài \texttt{dli\_sat} bằng công thức \((1 - \exp(-3 \cdot \text{DLI}/\text{dli\_sat})) / (1 - e^{-3})\); thừa số phạt sốc nhiệt tích lũy \(f_{\text{heat}} = \max(0,\; 1 - \text{heat\_penalty} \cdot \text{frac}_{T > T_{\max}})\) tính trên tỷ lệ số giờ thực tế vượt quá ngưỡng chịu nhiệt \(T_{\max}\); và thừa số chênh lệch nhiệt ngày -- đêm \(f_{\text{DIF}} \in [0{,}88;\; 1{,}02]\) phản ánh hiệu ứng \(\text{DIF} = T_{\text{ngày}} - T_{\text{đêm}}\) (tối ưu trong khoảng \(3{,}0\text{--}7{,}5\,^\circ\text{C}\)) lên tích lũy chất khô được ghi nhận từ thực nghiệm nhà kính WUR~\cite{wurChallengeDataset}. Cuối cùng, hệ số tiềm ẩn cố định theo từng nhà kính \(\text{gh\_effect} \sim \mathcal{N}(1{,}0,\; 0{,}05^2)\) (chặn trong \([0{,}82;\; 1{,}18]\)) đại diện cho sai khác kết cấu/canh tác riêng không đo được, kết hợp nhiễu sinh học phân phối chuẩn logarit \(\varepsilon \sim \text{LogNormal}(0,\; 0{,}1^2)\) đảm bảo năng suất luôn không âm trước khi quy đổi về đơn vị \(\text{kg/m}^2\) bằng mức chuẩn \texttt{reference\_kg\_m2} của từng loài cây.

    \item \textbf{Chuẩn mực thiết kế dữ liệu học máy chống rò rỉ và đối chiếu dữ liệu thực nghiệm WUR:} Để mô hình học được đúng quy luật sinh học thay vì ``học vẹt'' đặc điểm thống kê giả tạo, mỗi nhà kính giả lập xáo trộn và trồng đầy đủ cả 5 giống cây trên mọi vùng khí hậu (loại bỏ hoàn toàn rò rỉ tương quan giả giữa vùng miền và loại cây). Chuỗi dữ liệu theo giờ đi qua chính hàm \texttt{compute\_climate\_features} dùng chung nhằm đảm bảo khớp 100\% với lúc suy luận thực tế, đồng thời bảo tồn tính tất định hoàn toàn theo tham số \texttt{--seed} (mặc định \texttt{42}, sinh \(200\) nhà kính \(\times\; 5\) vụ \(= 1.000\) mẫu). Các mô-đun huấn luyện \texttt{ml/} và suy luận \texttt{worker/} tuyệt đối không được phép truy cập \texttt{growth.py} hay \texttt{crops\_growth.json}, coi bộ sinh dữ liệu như hộp đen thực tế; riêng trường phụ \texttt{meta.oracle\_relative\_yields} (tính năng suất kỳ vọng không nhiễu ở độ ẩm đất tối ưu cho cả 5 cây) chỉ được lưu phục vụ đo chất lượng xếp hạng (\texttt{Top-1 Accuracy}, \texttt{NDCG@3}). Nhờ gắn cố định \texttt{gh\_effect} theo \texttt{greenhouse\_key}, phép chia \texttt{GroupShuffleSplit} đánh giá trung thực khả năng tổng quát hóa sang những nhà kính chưa từng thấy. Song song đó, nhóm tích hợp thêm dữ liệu thực nghiệm công khai từ 22 buồng trồng qua 4 kỳ \textit{Autonomous Greenhouse Challenge} của Đại học Wageningen (WUR, Hà Lan)~\cite{wurChallengeDataset} thông qua mô-đun chuyển đổi \texttt{ml/wur\_adapter.py} (quy đổi \texttt{HumDef} \(\rightarrow\) \texttt{RH}, \texttt{PARin/Iglob} \(\rightarrow\) \texttt{lux}, \texttt{WC} \(\rightarrow\) \texttt{soil\_moisture}, chuẩn hóa theo buồng \texttt{Reference} và trích xuất qua cùng đường ống \texttt{climate\_features.py}) nhằm bổ sung tính thực tế cho tập huấn luyện và kiểm thử.
\end{enumerate}

\section{Hạ tầng triển khai và công cụ đảm bảo chất lượng}
\label{sec:c2-devops}

Toàn bộ cụm máy chủ chạy trên hệ điều hành \textbf{Ubuntu 24.04 LTS} được điều phối bằng \textbf{Docker Compose v2}~\cite{dockerDocs} với bản dựng nhiều chặng (\textit{multi-stage build}) từ \texttt{node:22-alpine} sang \textbf{Caddy 2.8-alpine}~\cite{caddyDocs}. Quyết định chọn Caddy làm điểm kết thúc HTTPS duy nhất nhằm bảo mật luồng xác thực JWT và kết nối WebSocket thời gian thực (\texttt{wss://}): Caddy tự động đăng ký và gia hạn chứng chỉ TLS Let's Encrypt, đồng thời phục vụ tệp tĩnh của Web Console và đóng vai trò reverse proxy chuyển tiếp \texttt{/api/*} cùng \texttt{/ws*} vào tiến trình \texttt{api}. Nhờ gắn volume \texttt{caddy\_data} ở chế độ chỉ đọc vào container Mosquitto kèm lịch cron khởi động lại Mosquitto mỗi tháng một lần, cả trình duyệt người dùng lẫn Gateway ESP32 đều dùng chung một chuỗi tin cậy chuẩn \texttt{ISRG Root X1} mà không phải duy trì hai hệ thống cấp phát chứng chỉ song song.

Trong suốt quá trình phát triển, nhóm duy trì kiểm soát chất lượng mã nguồn bằng \textbf{pytest} cùng \texttt{pytest-asyncio} (cấu hình tại \texttt{pyproject.toml}) cho 58 ca kiểm thử đơn vị và kiểm thử tích hợp end-to-end~\cite{pytestDocs}, kết hợp \textbf{ruff} để rà soát lỗi tĩnh Python và trình biên dịch \texttt{tsc} (cấu hình qua \texttt{web/tsconfig.json}) để kiểm tra chặt kiểu dữ liệu TypeScript ở phía giao diện.

\section{Tổng kết chương}
\label{sec:c2-summary}

Trong Chương 2, nhóm chúng em đã trình bày bức tranh tổng thể về giải pháp kỹ thuật 4 tầng, bộ công nghệ thực tế kèm phân tích các quyết định kiến trúc gắn liền với từng phân hệ của GreenHub. Từ việc chọn ESP-NOW ở tầng thiết bị biên nhằm tối ưu năng lượng và gom an toàn về Gateway, qua tầng kết nối Eclipse Mosquitto TLS dùng Dynamic Security ACL chia sẻ chứng chỉ từ Caddy, đến kiến trúc backend monolith hai tiến trình tách biệt RAM cùng PostgreSQL/TimescaleDB \texttt{LISTEN/NOTIFY}, giao diện React 18 PWA cập nhật thời gian thực qua WebSocket, và phân hệ LightGBM với mô-đun đặc trưng Python thuần đi kèm bộ sinh dữ liệu nông học chuẩn hóa đa nguồn, mỗi lựa chọn đều hướng tới độ tin cậy, tính nhất quán dữ liệu và khả năng vận hành tiết kiệm trên VPS. Nền tảng công nghệ này sẽ là cơ sở để nhóm triển khai phân tích yêu cầu và thiết kế chi tiết từng phân hệ trong Chương 3.